% =====================================================================
%  MATH300 Project Description -- Incremental SVD
% =====================================================================
\documentclass[11pt,a4paper]{article}

\usepackage[T1]{fontenc}
\usepackage[utf8]{inputenc}
\usepackage{lmodern}
\usepackage[margin=2.5cm]{geometry}
\usepackage{enumitem}
\usepackage{booktabs}
\usepackage{tabularx}
\usepackage{array}
\usepackage{amsmath,amssymb}
\usepackage{xcolor}
\usepackage[colorlinks=true,urlcolor=blue,linkcolor=black,citecolor=black]{hyperref}

\setlist{itemsep=2pt, topsep=4pt}
\setlength{\parindent}{0pt}
\setlength{\parskip}{6pt}

% Placeholder macro: blue text in square brackets.
% To hide the colour in the final version, change {blue} to {black}.
\newcommand{\ph}[1]{\textcolor{blue}{[#1]}}

% Marker for optional bonus tasks.
\newcommand{\bonus}{\textbf{(Bonus)}}

\newcolumntype{L}{>{\raggedright\arraybackslash}X}

\begin{document}

% ---------------------------------------------------------------------
%  Title (no logo, no instructor name)
% ---------------------------------------------------------------------
\begin{center}
  {\large\bfseries MATH300: Mathematical Computing \textperiodcentered{} Fall 2026}\\[6pt]
  {\Large\bfseries Project: Compressing a Video with the (Incremental) SVD}
\end{center}
\vspace{6pt}

Groups must be formed, and project titles submitted on Canvas, by
\textbf{October 26, 2026}. The project
has two graded deliverables and an oral exam: a 3--5 page report for
Project~2 (due \textbf{November 20, 2026}), a poster for the final project
(due \textbf{December 11, 2026}), and an oral exam held between November 30
and December 17, 2026, 3:30\,pm. Read the general rules below before you
start; they apply throughout the project.

% ---------------------------------------------------------------------
\section*{General remarks}
% ---------------------------------------------------------------------
\begin{itemize}
  \item Projects are set for a group of \textbf{two or three} students.
  \item Form your group by \textbf{October 26, 2026}, and submit your
        project title as a group assignment on \textbf{Canvas} by the same
        date.
  \item Every group member must carry an \textbf{equal portion} of the coding
        and of the report and poster preparation.
  \item The scope can be adjusted to the number of participants, their
        interests and prior knowledge.
  \item Only \textbf{Python} may be used for programming.
  \item Implementation, documentation, validation and demonstration examples
        are all part of the project.
  \item Within these rules you are free to choose the implementation, code
        structure and exact content. You can discuss this with me at any time.
\end{itemize}

% ---------------------------------------------------------------------
\section*{Version control and GitHub}
% ---------------------------------------------------------------------
Each group keeps \textbf{one GitHub repository for the whole project}
(Project~2 and the final project), and every member commits from their own
GitHub account.
\begin{itemize}
  \item Use git from the very beginning. Commits should be small, consistent
        and clearly described. Avoid large commits.
  \item Each group member uploads their own code using \textbf{their own
        GitHub account}. Commits from a shared or borrowed account will not
        be credited to the other members.
  \item The repository name is up to your group.
  \item If the repository is private, add my GitHub account
        \textbf{\texttt{say-yas}} as a collaborator.
  \item The repository must contain a README describing the project, how to
        install and run the code, and the AI-use statement (see below).
\end{itemize}

% ---------------------------------------------------------------------
\section*{Use of AI tools}
% ---------------------------------------------------------------------
\textbf{AI-generated code is not accepted in any form.} Code produced by a
chat assistant or code generator (e.g.\ ChatGPT, Claude, Gemini) and copied,
pasted or adapted into your project will be treated as work that is not your
own.
\begin{itemize}
  \item \textbf{Permitted:} inline coding assistance such as GitHub Copilot
        autocompletion while you write the code yourself, \textbf{provided its
        use is acknowledged}.
  \item \textbf{Acknowledgement:} include an AI-use statement in the README
        and in the report that names the tool(s) used and the parts of the
        code (files or functions) where they assisted. If no AI tools were
        used, state this explicitly.
  \item \textbf{Reports and posters must not be generated by AI} in any
        form: their text, figures and layout must be your own work. If AI
        tools were used for any other non-code purpose (e.g.\ checking
        grammar or searching the literature), describe this in the same
        statement.
  \item You are responsible for every line in your repository. In the oral
        exam you must be able to explain and modify any part of the code you
        submitted.
\end{itemize}

% ---------------------------------------------------------------------
\section*{Deliverables and deadlines}
% ---------------------------------------------------------------------
Submit each deliverable following the details provided on the
\textbf{Canvas} page. The written report and poster are handed in on
\textbf{Crowdmark}, and the link to your GitHub repository is submitted
separately on \textbf{Canvas}.

\begin{center}
\small
\renewcommand{\arraystretch}{1.35}
\begin{tabularx}{\textwidth}{|>{\raggedright\arraybackslash}p{2.5cm}|>{\raggedright\arraybackslash}p{3.0cm}|>{\raggedright\arraybackslash}p{3.0cm}|L|}
\hline
\textbf{Deliverable} & \textbf{Due} & \textbf{Submission} & \textbf{Notes} \\
\hline\hline
Group and project title & Monday, October 26, 2026 &
  Group assignment on Canvas &
  Group of two or three members; submit the project title. \\
\hline
Project~2 report & Friday, November 20, 2026 &
  Report (PDF) on Crowdmark; repository link on Canvas &
  3--5 pages; the page limit \textbf{includes} references, figures and
  appendices. \\
\hline
Final project poster & Friday, December 11, 2026 &
  Poster (PDF) on Crowdmark; repository link on Canvas &
  36 inches high by 48 inches wide (landscape). Target audience:
  undergraduate classmates with a basic background in math, physics or
  engineering. \\
\hline
Oral exam & Between November 30 and December 17, 2026, 3:30\,pm &
  Individual meeting; details discussed in class &
  30 minutes per student. Your poster must already be submitted on
  Crowdmark before your exam. The last possible slot ends December 17 at
  3:30\,pm. \\
\hline
\end{tabularx}
\end{center}

\subsection*{Oral exam}
The oral exam is individual and lasts 30 minutes per student. It takes place
between November 30 and December 17, 2026, 3:30\,pm. The oral exams and how
to arrange them will be discussed in class before the Thanksgiving break.

Before your oral exam, your poster must also be submitted on
\textbf{Crowdmark}; that is, you must have finished both Project~2 and the
final project. In the meeting you will:
\begin{itemize}
  \item Present your results and algorithms.
  \item Possibly live-code and run parts of your algorithms to generate
        outputs on the spot.
  \item Answer questions about the theory, your implementation and your
        results. Each member should be able to explain the whole project, not
        only their own parts.
\end{itemize}
Bring a laptop with the repository cloned and a working environment, so that
the code runs without additional setup.

% ---------------------------------------------------------------------
\section*{Grading}
% ---------------------------------------------------------------------
\begin{center}
\small
\renewcommand{\arraystretch}{1.35}
\begin{tabularx}{0.8\textwidth}{|L|>{\raggedleft\arraybackslash}p{3.2cm}|}
\hline
\textbf{Component} & \textbf{Points} \\
\hline\hline
Group and project title (on time) & 10 \\
\hline
Project~2 report & 100 \\
\hline
Final project poster & 100 \\
\hline
Oral exam (individual) & 50 \\
\hline
Bonus: extra steps beyond the required tasks & up to 10 \\
\hline
\end{tabularx}
\end{center}

The Project~2 report and the poster are each graded out of 100. What matters
most is:
\begin{itemize}
  \item \textbf{Coherence:} a clear story from the question, through the
        methods, to the results and conclusions.
  \item \textbf{Correctness:} of the mathematics, the code and the results.
  \item \textbf{Depth of analysis and discussion:} explain what your results
        mean, why they look the way they do, and what their limitations are,
        rather than only showing them.
  \item \textbf{Professional presentation:} figures with labeled axes and
        readable font sizes, captions and legends, a clean and readable
        layout, and proper references.
\end{itemize}

A complete, correct and well-discussed
solution of the required tasks can earn full marks. Tasks marked
\bonus{} and the suggestions under ``Going further'' are optional; exploring
such extra steps can earn \textbf{up to 10 bonus points} in total.

% ---------------------------------------------------------------------
\section*{Requirements for the report and the poster}
% ---------------------------------------------------------------------
Both the report and the poster must contain a clear description of the
project, references, properly labeled figures and an author-contributions
statement.

\subsection*{Content}
\begin{itemize}
  \item A clear and concise description of the problem, the methods and the
        results.
  \item The AI-use statement described above.
  \item The report should not contain any significant amount of code; refer
        to the repository instead.
\end{itemize}

\subsection*{Poster}
\begin{itemize}
  \item The poster must be \textbf{self-contained}: a reader should be able to
        follow it without the report, the code or a spoken explanation.
  \item You may be selective in what you show and how you narrate your story,
        so that an undergraduate audience with a basic background in math,
        physics or engineering can follow it.
  \item You may choose your own title for the poster; be creative.
\end{itemize}

\subsection*{References}
\begin{itemize}
  \item Include a list of references at the end of the report and the poster.
  \item Cite sources clearly where they are used in the text (papers, books,
        software, code repositories and data sets).
  \item Use one consistent citation style throughout.
\end{itemize}

\subsection*{Figures}
\begin{itemize}
  \item Every plot must have labeled $x$ and $y$ axes (and $z$ axis for 3D
        plots or color bars), including units where applicable.
  \item Axis labels, tick labels and legends must have a font size similar to
        the surrounding text, so they remain readable without zooming.
  \item Each figure needs a legend when it shows more than one data series,
        and a caption that states what is shown.
  \item Use log scales where appropriate (e.g.\ for convergence plots) and say
        so in the caption.
\end{itemize}

\subsection*{Author contributions}
Each group must include an author-contributions statement in the report and
the README that states:
\begin{itemize}
  \item Who did what: which parts of the code, experiments, analysis, writing
        and poster each member worked on.
  \item To what extent: the approximate share of each member in each part
        (e.g.\ 50\,\% / 50\,\%).
  \item Who took the lead in which parts.
\end{itemize}
The coding and the report and poster preparation must be shared equally
between the members. The commit history should be consistent with the stated
contributions.

% ---------------------------------------------------------------------
\section*{Codebase requirements}
% ---------------------------------------------------------------------
\begin{itemize}
  \item Well structured and documented, both in the code (docstrings,
        comments) and in the repository (README).
  \item Validated: show that your implementation produces correct results,
        e.g.\ on examples with known solutions.\end{itemize}

% ---------------------------------------------------------------------
\section*{Project description}
% ---------------------------------------------------------------------

\subsection*{Goal}
A video is a sequence of frames, and consecutive frames are usually very
similar. In this project you will use the singular value decomposition (SVD)
to compress a short video: you will approximate all its frames by a few
basis images. First you compute the SVD of the whole video at once with
NumPy; then you implement the \emph{incremental SVD} described in
Section~2.2 of Kühl et al.\ \cite{kuehl}, which processes the frames a few at
a time, as if they arrived one after another, and never stores the whole
video. Alternatively, you can track the SVD of the frames themselves, one
frame after another, by treating the change between consecutive frames as a
small modification of the previous frame (Approach~B below). You do
\textbf{not} need to implement the SVD itself: use
\texttt{numpy.linalg.svd} \cite{numpy} wherever an SVD is needed.

\subsection*{Background}
Convert each frame to grayscale and reshape it into one long column vector
of $N$ pixel values. Placing the $T$ frames side by side gives the matrix
$Y \in \mathbb{R}^{N \times T}$, one column per frame, with usually
$N \gg T$. Its (reduced) SVD is
\begin{equation*}
  Y = U S V^{\mathsf{T}},
\end{equation*}
where $U$ and $V$ have orthonormal columns and $S$ is diagonal with the
singular values $s_1 \ge s_2 \ge \dots \ge 0$. The columns of $U$ are basis
images, and each frame is a combination of them. Keeping only the $q$ largest
singular values and the corresponding columns of $U$ and $V$ gives the best
approximation of $Y$ of rank $q$. If the singular values decay quickly, a
small $q$ already reproduces the video well, and storing the truncated $U$,
$S$ and $V$ needs much less memory than storing all frames.

For a long video, $Y$ may not fit in memory. The incremental SVD
\cite{kuehl,brand2002,brand2006} instead builds the approximation step by
step. You may choose one of the following two approaches.

\textbf{Approach A: adding frames as new columns.} Suppose we already have $U$, $S$ and $V$ for the frames seen so far,
and a new block $B \in \mathbb{R}^{N \times b}$ of $b$ frames arrives.
Following \cite[Section 2.2]{kuehl}:
\begin{enumerate}
  \item Project the new frames onto the current basis: $M = U^{\mathsf{T}} B$.
  \item Compute the part that is new, $P = B - U M$, and its QR
        decomposition $P = Q_P R_P$.
  \item Form the small matrix
        \begin{equation*}
          K = \begin{bmatrix} S & M \\ 0 & R_P \end{bmatrix}
        \end{equation*}
        and compute its SVD, $K = U' S' V'^{\mathsf{T}}$.
  \item Update the factors:
        \begin{equation*}
          U \leftarrow [\,U \;\; Q_P\,]\, U', \qquad
          S \leftarrow S', \qquad
          V \leftarrow \begin{bmatrix} V & 0 \\ 0 & I \end{bmatrix} V'.
        \end{equation*}
  \item Truncate: keep only the $q$ largest singular values and the
        corresponding columns of $U$ and $V$.
\end{enumerate}
Only the small matrix $K$ is decomposed, so each update is cheap. In exact
arithmetic, and without truncation, the result is the same as the SVD of all
frames at once. In floating-point arithmetic the columns of $U$ slowly lose
their orthogonality, especially when $P$ is nearly zero, so in practice one
re-orthogonalizes $U$ from time to time \cite{kuehl}.

\textbf{Approach B: updating one frame to the next.} Here each frame is kept
as an image matrix $F_k$ (pixels in rows and columns). Start from the SVD of
the first frame, $F_1 = U S V^{\mathsf{T}}$. The next frame is
$F_{k+1} = F_k + D_k$, where the difference $D_k = F_{k+1} - F_k$ is usually
small. Write (or approximate) the difference as a product
$D_k \approx B C^{\mathsf{T}}$ with only $b$ columns in $B$ and $C$, for
example from the $b$ largest singular values of $D_k$. Then update the SVD
with the general formula for $Y \to Y + B C^{\mathsf{T}}$ in
\cite[Section 2.2]{kuehl}:
\begin{enumerate}
  \item Compute $M_B = U^{\mathsf{T}} B$ and the QR decomposition
        $B - U M_B = Q_B R_B$, and likewise $M_C = V^{\mathsf{T}} C$ and
        $C - V M_C = Q_C R_C$.
  \item Form the small matrix
        \begin{equation*}
          K = \begin{bmatrix} S & 0 \\ 0 & 0 \end{bmatrix}
            + \begin{bmatrix} M_B \\ R_B \end{bmatrix}
              \begin{bmatrix} M_C \\ R_C \end{bmatrix}^{\mathsf{T}}
        \end{equation*}
        and compute its SVD, $K = U' S' V'^{\mathsf{T}}$.
  \item Update: $U \leftarrow [\,U \;\; Q_B\,]\, U'$, $S \leftarrow S'$,
        $V \leftarrow [\,V \;\; Q_C\,]\, V'$, and truncate to rank $q$.
\end{enumerate}
This gives a low-rank approximation of every frame without computing a new
SVD of each frame. With truncation, the small errors made in each step can
add up from frame to frame.

\subsection*{Tasks for Project 2 (report, due November 20, 2026)}
\begin{enumerate}
  \item Choose a short video and use at most 100 frames. Read its
        frames into Python (for example with the \texttt{imageio} package),
        convert them to grayscale, make them small (for example
        $100 \times 100$ pixels), and build the matrix $Y$.
  \item Compute the SVD of $Y$ with \texttt{numpy.linalg.svd} \cite{numpy}
        and plot the singular values (log scale).
  \item Implement the incremental SVD with Approach~A or Approach~B above,
        with optional truncation to rank $q$. In Approach~A, process the
        frames in blocks of $b$ frames (for example $b = 10$). Use
        \texttt{numpy.linalg.svd} and \texttt{numpy.linalg.qr} for the small
        decompositions.
  \item Run it \emph{without} truncation and check your result against
        \texttt{numpy.linalg.svd}: in Approach~A, compare the singular
        values with those of the full SVD of $Y$; in Approach~B, compare the
        singular values of a few frames with the SVD of those frames. (If
        they differ, the columns of $U$ have probably lost their
        orthogonality; see the remark above.)
  \item \bonus\ Check how orthogonal the columns of $U$ stay by computing
        $\lVert U^{\mathsf{T}} U - I \rVert_F$ after each update, and add a
        re-orthogonalization step if needed.
  \item \bonus\ Implement both approaches.
\end{enumerate}

\subsection*{Tasks for the final project (poster, due December 11, 2026)}
\begin{enumerate}
  \item Compare rank-$q$ approximations from the incremental SVD (with
        truncation) with those from \texttt{numpy.linalg.svd}, for a few
        values of $q$, for example $q = 1, 5, 20$. Show one original frame
        next to its reconstructions, and plot the relative error against $q$
        (Approach~A: $\lVert Y - Y_q \rVert_F / \lVert Y \rVert_F$ for the
        whole video; Approach~B: the error of each frame, also against the
        frame number).
  \item Discuss what you observe: is the incremental SVD as good as the full
        SVD? What are its benefits (for example memory) and shortcomings?
  \item \bonus\ Study how the block size $b$ affects the accuracy and the
        computing time.
\end{enumerate}

\subsection*{Going further}
These optional steps count toward the bonus points. For example, you can show
the first few columns of $U$ as images and explain what they represent.

% ---------------------------------------------------------------------
\begin{thebibliography}{9}
  \bibitem{kuehl} N.~Kühl, H.~Fischer, M.~Hinze and T.~Rung (2023). An
    incremental singular value decomposition approach for large-scale
    spatially parallel \& distributed but temporally serial data -- applied
    to technical flows. arXiv:2302.09149.
    \url{https://arxiv.org/abs/2302.09149}
  \bibitem{brand2002} M.~Brand (2002). Incremental singular value
    decomposition of uncertain data with missing values. In
    \emph{Computer Vision -- ECCV 2002}, Lecture Notes in Computer Science,
    vol.~2350, Springer. \url{https://doi.org/10.1007/3-540-47969-4_47}
  \bibitem{brand2006} M.~Brand (2006). Fast low-rank modifications of the
    thin singular value decomposition. \emph{Linear Algebra and its
    Applications}, 415, 20--30.
    \url{https://www.merl.com/publications/TR2006-059}
  \bibitem{numpy} NumPy. \emph{numpy.linalg.svd}.
    \url{https://numpy.org/doc/stable/reference/generated/numpy.linalg.svd.html}
\end{thebibliography}

\end{document}